
\documentclass{article}
\usepackage{colortbl}
\usepackage{makecell}
\usepackage{multirow}
\usepackage{supertabular}

\begin{document}

\newcounter{utterance}

\centering \large Interaction Transcript for game `eqbench', experiment `eq\_bench', episode 151 with reflexion{-}llama{-}3.1{-}8b{-}instruct.
\vspace{24pt}

{ \footnotesize  \setcounter{utterance}{1}
\setlength{\tabcolsep}{0pt}
\begin{supertabular}{c@{$\;$}|p{.15\linewidth}@{}p{.15\linewidth}p{.15\linewidth}p{.15\linewidth}p{.15\linewidth}p{.15\linewidth}}
    \# & \multicolumn{2}{c}{Player} && \multicolumn{2}{c}{Game Master} \\
    \hline

    \theutterance \stepcounter{utterance}  
    & & & \multicolumn{4}{p{0.6\linewidth}}{
        \cellcolor[rgb]{0.9,0.9,0.9}{
            \makecell[{{p{\linewidth}}}]{
                \texttt{\tiny{[P1$\langle$GM]}}
                \texttt{You have attempted this same task before and did not succeed. These are the plans you wrote after those attempts.} \\
\texttt{Trial \#1: New plan: } \\
\\ 
\texttt{Upon re{-}reading the dialogue, I realized that I focused too much on Ellie's reactions and didn't consider Dr. Rigby's perspective enough. I should have paid more attention to the clues in the dialogue that indicate Dr. Rigby's emotional state. Specifically, Dr. Rigby's tone and language suggest a sense of conviction and passion, which is more consistent with being Eager rather than Defensive. Additionally, Dr. Rigby's use of metaphors and hypothetical scenarios implies a sense of excitement and enthusiasm, which further supports the Eager emotion. } \\
\\ 
\texttt{To improve my strategy, I will focus on analyzing the speaker's tone, language, and behavior to better understand their emotional state. I will also consider the context of the dialogue and the relationships between the characters to gain a deeper understanding of the emotional dynamics at play.} \\
\\ 
\\ 
\\ 
\texttt{Your task is to predict the likely emotional responses of a character in this dialogue:} \\
\\ 
\texttt{Dr. Rigby: You remind me of someone I used to know.} \\
\texttt{Ellie: Oh, yeah? Some loser you used to pity?} \\
\texttt{Dr. Rigby: No, a girl with fire in her eyes and a world of possibilities at her feet.} \\
\texttt{Ellie: Sounds like a fairy tale. No one's got a world of possibilities.} \\
\texttt{Dr. Rigby: Not with that attitude, no. But imagine, just for a moment, a different reality.} \\
\texttt{Ellie: I don't have time for your riddles, lady.} \\
\texttt{Dr. Rigby: It's not a riddle, it's a wormhole.} \\
\texttt{Ellie: A what?} \\
\texttt{Dr. Rigby: A wormhole. A shortcut through space and time. What if you could reach out and change your destiny?} \\
\texttt{Ellie: You're nuts.} \\
\texttt{Dr. Rigby: Perhaps. But what if I told you that the future is not set in stone? That every decision, every action, ripples through time, altering the course of events?} \\
\texttt{Ellie: And what? I'm supposed to believe that my shitty decisions matter?} \\
\texttt{Dr. Rigby: More than you can imagine. You have the power to shape your own future.} \\
\texttt{Ellie: And how would you know?} \\
\texttt{{[}End dialogue{]}} \\
\\ 
\texttt{At the end of this dialogue, Dr. Rigby would feel...} \\
\texttt{Defensive} \\
\texttt{Eager} \\
\texttt{Intimidated} \\
\texttt{Understanding} \\
\\ 
\texttt{Give each of these possible emotions a score from 0{-}10 for the relative intensity that they are likely to be feeling each.} \\
\\ 
\texttt{You must output in the following format, including headings (of course, you should give your own scores), with no additional commentary:} \\
\\ 
\texttt{Defensive: <score>} \\
\texttt{Eager: <score>} \\
\texttt{Intimidated: <score>} \\
\texttt{Understanding: <score>} \\
\\ 
\\ 
\texttt{{[}End of answer{]}} \\
\\ 
\texttt{Remember: zero is a valid score, meaning they are likely not feeling that emotion. You must score at least one emotion > 0.} \\
\\ 
\texttt{Your answer:} \\
            }
        }
    }
     \\ \\

    \theutterance \stepcounter{utterance}  
    & \multicolumn{4}{p{0.6\linewidth}}{
        \cellcolor[rgb]{0.561,0.737,0.561}{
            \makecell[{{p{\linewidth}}}]{
                \texttt{\tiny{[P1$\rangle$GM]}}
                \texttt{Defensive: 0} \\
\texttt{Eager: 9} \\
\texttt{Intimidated: 0} \\
\texttt{Understanding: 2} \\
            }
        }
    }
    & & \\ \\

    \theutterance \stepcounter{utterance}  
    & & & \multicolumn{2}{p{0.3\linewidth}}{
        \cellcolor[rgb]{0.9,0.9,0.9}{
            \makecell[{{p{\linewidth}}}]{
                \texttt{\tiny{[GM$|$GM]}}
                \texttt{Defensive: 0} \\
\texttt{Eager: 9} \\
\texttt{Intimidated: 0} \\
\texttt{Understanding: 2} \\
            }
        }
    }
    & & \\ \\

    \theutterance \stepcounter{utterance}  
    & & & \multicolumn{2}{p{0.3\linewidth}}{
        \cellcolor[rgb]{0.9,0.9,0.9}{
            \makecell[{{p{\linewidth}}}]{
                \texttt{\tiny{[GM$|$GM]}}
                \texttt{Defensive: 0} \\
\texttt{Eager: 7} \\
\texttt{Intimidated: 0} \\
\texttt{Understanding: 8} \\
            }
        }
    }
    & & \\ \\

    \theutterance \stepcounter{utterance}  
    & & & \multicolumn{2}{p{0.3\linewidth}}{
        \cellcolor[rgb]{0.9,0.9,0.9}{
            \makecell[{{p{\linewidth}}}]{
                \texttt{\tiny{[GM$|$GM]}}
                \texttt{game\_result = LOSE} \\
            }
        }
    }
    & & \\ \\

    \theutterance \stepcounter{utterance}  
    & & \multicolumn{4}{p{0.6\linewidth}}{
        \cellcolor[rgb]{0.95,0.95,0.95}{
            \makecell[{{p{\linewidth}}}]{
                \texttt{\tiny{[P1$\leftrightarrow$P]}}
                \texttt{Post{-}episode Reflexion (not sent to GM):} \\
\texttt{New plan: } \\
\\ 
\texttt{Upon re{-}reading the dialogue, I realized that I focused too much on Ellie's reactions and didn't consider Dr. Rigby's perspective enough. I should have paid more attention to the clues in the dialogue that indicate Dr. Rigby's emotional state. Specifically, Dr. Rigby's tone and language suggest a sense of conviction and passion, which is more consistent with being Eager rather than Defensive. Additionally, Dr. Rigby's use of metaphors and hypothetical scenarios implies a sense of excitement and enthusiasm, which further supports the Eager emotion. } \\
\\ 
\texttt{To improve my strategy, I will focus on analyzing the speaker's tone, language, and behavior to better understand their emotional state. I will also consider the context of the dialogue and the relationships between the characters to gain a deeper understanding of the emotional dynamics at play.} \\
\\ 
\texttt{To better score the emotions, I will consider the following:} \\
\\ 
\texttt{{-} Eager: I will look for language and tone that suggests enthusiasm, conviction, and a sense of excitement. Dr. Rigby's use of metaphors and hypothetical scenarios, such as the "wormhole" and "shortcut through space and time," suggests a sense of passion and conviction.} \\
\texttt{{-} Defensive: I will look for language and tone that suggests defensiveness, such as a rise in pitch, a change in tone, or a shift in language that implies a sense of being attacked or criticized. However, in this dialogue, Dr. Rigby's tone and language do not suggest defensiveness.} \\
\texttt{{-} Intimidated: I will look for language and tone that suggests fear, anxiety, or a sense of being overwhelmed. However, in this dialogue, Dr. Rigby's tone and language do not suggest intimidation.} \\
\texttt{{-} Understanding: I will look for language and tone that suggests empathy, compassion, or a sense of connection. While Dr. Rigby's language and tone do not suggest a deep understanding of Ellie's perspective, they do suggest a sense of conviction and passion, which is more consistent with being Eager.} \\
\\ 
\texttt{With this new plan, I will re{-}score the emotions as follows:} \\
\\ 
\texttt{Defensive: 0} \\
\texttt{Eager: 10} \\
\texttt{Intimidated: 0} \\
\texttt{Understanding: 4} \\
            }
        }
    }
    & \\ \\

\end{supertabular}
}

\end{document}
